%% ======================================================================

%% -------------------------
%% Theme
%% -------------------------
\documentclass[aspectratio=169]{beamer}
\usetheme{Madrid}
\definecolor{hmt-red}{RGB}{178,42,46}
\colorlet{beamer@blendedblue}{hmt-red}
\newcommand{\bu}[1]{\textbf{\uline{#1}}}

%% -------------------------
%% Packages
%% -------------------------
\usepackage{graphicx}
\usepackage{booktabs}
\usepackage{amsmath,amssymb}
\usepackage{adjustbox}
\usepackage{multicol}
\usepackage{ulem}
\usepackage{enumitem}
\usepackage{hyperref}

\hypersetup{
    colorlinks=false,           % Must be false to show the border/underline
    linkcolor=hmt-red,          % Color for internal links
    urlcolor=hmt-red,           % Color for external URLs
    pdfborder={0 0 1},          % Sets border width: {horizontal vertical thickness}
    pdfborderstyle={/S/U/W 1},  % /S/U sets style to Underline; /W 1 sets width to 1pt
    allbordercolors=hmt-red     % Applies the red color to the underline
}

%% -------------------------
%% Bibliography
%% -------------------------
\usepackage[
  backend=biber,
  style=apa,
]{biblatex}
\addbibresource{references.bib}

%% -------------------------
%% Metadata (EDIT HERE ONCE)
%% -------------------------

\title
    [Agglomeration \& rent-sharing]
    {Spatial misallocation under agglomeration externalities, isolating the causal effect of rent-sharing}
\subtitle{MSc Economics dissertation proposal}
\author[Aloysius Lip]
    {Aloysius Lip\inst{1,2}         \\	[1ex]
    	{Supervisor: Lars Nesheim\inst{1,3}}}
\institute[UCL] {
    \inst{1}
    Department of Economics,
    University College London (UCL)
    \and
    \inst{2}
    HM Treasury (HMT)
    \and
    \inst{3}
    Centre for Microdata Methods and Practice (cenmap)
}
\date{February 2026}
% \titlegraphic{\includegraphics[width=1em]
%     {ucl_logo.jpg}
% }

%------%

\begin{document}
\maketitle
\section{Introduction} 

\begin{frame}{Overview}
	\begin{itemize}
		\item \textbf{Core}: an ML exercise process to estimate varied distance and firm characteristic effects on agglomeration.
		We want to measure strength and heterogeneity.
		\item \textbf{Analysis}: we need valid causal inference, so draw on the following fields of research
			\begin{itemize}
				\item \bu{Spatial/urban econometrics} (applying learning from \bu{Time series econometrics} course) as appropriate
				\item \bu{Macroeconomic policy} research learnt for identification of exogenous shocks
			\end{itemize}
		\item \textbf{Extensions}:
			\begin{enumerate}[label=(\alph*)]
				\item Clustering regressor (\bu{ML in economics})
				\item Regularisation to remove ineffective characteristic interactions from neighbouring firms (\bu{ML})
				\item Literature from social peer effects literature (\bu{Labour economics})
				\item Introduce regressor: exogenously-determined share of landlord bargaining power
			\end{enumerate}
	\end{itemize}
\end{frame}

\begin{frame}{Hypothesis}
	\begin{itemize}
		\item Sector heterogeneity matters for agglomeration, but less than expected - \\
				dense firms of all kinds still share knowledge
		\item Other measures of distance beyond spatial matter for agglomeration - \\
				transport links are a good substitute
		\item There are multiple equilibria for agglomeration steady-states: 
		\item Landlord bargaining power is a key determinant - \\
				more landlord LVC means less agglomeration, and more misallocation
	\end{itemize}
\end{frame}

\begin{frame}{Motivation (academia)}
    \begin{block}{Field}
        \textbf{Macroeconomics: Growth \& Productivity}
        \begin{itemize}
            \item The role of agglomeration in the firm production function, including spillovers
            \item Drawing from \textbf{Urban Economics}
            \item Analysis using \textbf{Microeconometric} techniques
        \end{itemize}
    \end{block}
\end{frame}

\begin{frame}{Motivation (policy)}
    \begin{itemize}
        \item Local Growth Plans in 12 Combined Authorities in England$^1$
    \end{itemize}
    \begin{block}{HMT AoRI$^2$}
        \begin{itemize}
            \item \textbf{1.15}: What policies and interventions would be most effective in driving regional and local economic growth, and how do these vary across places?
            \item \textbf{3.2}: How can the financial sector best support the real economy, both locally in the places financial firms agglomerate, but also across the country?
					Are there ways in which the sector’s support for the real economy could be improved?
        \end{itemize}
    \end{block}
    \footnotetext[1]{\href{https://www.gov.uk/government/publications/english-devolution-white-paper-power-and-partnership-foundations-for-growth/english-devolution-white-paper\#local-growth-plans}{English Devolution White Paper}}
    \footnotetext[2]{\href{https://www.gov.uk/government/publications/hm-treasury-areas-of-research-interest}{HM Treasury: Areas of Research Interest}}
\end{frame}
%-----%

\begin{frame}{Macroeconomic facts}
    \begin{block}{Firm density vs productivity growth (cities, time)}
    
    \end{block}
    \begin{block}{Density vs productivity (global cities)}
    
    \end{block}
\end{frame}
%-----%

\begin{frame}{Literature}
\end{frame}
%-----%

\begin{frame}{Identification}
\end{frame}
%-----%

\begin{frame}{Models}
    \framehead{3-factor production function}
    \begin{align}
        Y_t &= A_t K_t^\alpha L_t^\beta Z_t^{\gamma}    \\
        Z_t(r) &= \int_0^S\psi(r,s)\cdot s \cdot \theta(s)n(s)\ ds
    \end{align}
    \begin{itemize}
        \item $Z_t(r)$: external agglomeration technology
    \end{itemize}
\end{frame}
%-----%

\begin{frame}{Reduced form estimation}

\end{frame}
%-----%

\begin{frame}{Applied teaching}
    \begin{block}{Macroeconomic Policy}
    
    \end{block}
    \begin{block}{Machine Learning in Economics}
    
    \end{block}
    \begin{block}{Time Series Econometrics}
    \end{block}
    \begin{block}{Labour Economics}
    	\begin{itemize}
    
    	\end{itemize}
    \end{block}
\end{frame}

\begin{frame}{Data (sources)}
    \begin{table}
	\centering
    \begin{tabular}{l p{3.5cm} c c r}
        \toprule
          \textbf{Source}
        & \textbf{Granularity}
        & \textbf{Location}
        & \textbf{Period}   
        & \textbf{Access}    \\
        
        \midrule

        1. FAME
        & Medium \& large firms
        & HQs only
        & 2006-2025 (20-year rolling)
        & Public
        \\

        2. ONS - ABS
        & Sample of 2000 firms
        & \checkmark
        & 
        & SRS (via DBT)
        \\

        3. ONS - all
        & All firms
        & \checkmark
        &
        & SRS (via DBT)
        \\

        4. HMT - travel times
        & 1km2 gridsquares
        & Non-firm specific
        & Instrument (1Y cross-sectional)
        & Non-public (via HMT)
        
		\bottomrule
    \end{tabular}
    \end{table}
\end{frame}
%-----%

\begin{frame}{Data (calibration \& estimation)}
\end{frame}
%-----%

\end{document}
